\section{Shielding consumers: cap versus transfer}
\label{sec:shield}

Both instruments in the recent literature stabilise the crisis, and
Table~\ref{tab:summary} shows they do so comparably. They part ways on the margin
this paper adds.

\subsection{The cap eliminates the adoption margin}
\label{sec:cap}

Under the cap the peak green-share response falls from $8.9$ pp to essentially
zero ($-0.04$ pp) under \emph{both} shocks: the margin is not attenuated, it is
switched off. The mechanism is direct---the cap holds $P^{E}_{B}$ at its
pre-crisis level, and $P^{E}_{B}$ is precisely what determines the return to
adoption.

\begin{figure}[H]\centering
\includegraphics[width=0.92\textwidth]{output/fig2_policy_price.png}
\caption{Fiscal response to the \emph{price} shock, adoption open: no policy
(solid), price cap (dashed), Slutsky transfer (dotted).}
\label{fig:policy_price}
\end{figure}

\begin{figure}[H]\centering
\includegraphics[width=0.92\textwidth]{output/fig2_policy_supply.png}
\caption{Fiscal response to the \emph{supply} shock, adoption open: no policy
(solid), price cap (dashed), Slutsky transfer (dotted).}
\label{fig:policy_supply}
\end{figure}

%\begin{figure}[H]\centering
%\includegraphics[width=\textwidth]{output/fig3_policy_adoption.png}
%\caption{Policy and the adoption margin. Under both shocks the price cap
%(dashed) flattens the green-share response to zero; the transfer (dotted) leaves
%it intact.}
%\label{fig:policy_adoption}
%\end{figure}

\subsection{The transfer preserves it}
\label{sec:transfer}

The Slutsky transfer leaves the green-share response essentially unchanged
($9.1$ pp against $8.9$ with no policy) while delivering more output
stabilisation than the cap at similar fiscal cost ($\sum y=-22.3$ vs.\ $-53.0$;
gross fiscal $49.3$ vs.\ $53.9$). This is a direct consequence of the transfer's
lump-sum indexing (Section~\ref{sec:bop}): because it does not condition on the
household's current durable, switching forfeits nothing and the relative-price
incentive to adopt is left intact. Subsidies destroy the transition margin;
transfers do not.

\subsection{The cap is a dial, not a switch}
\label{sec:dial}

Table~\ref{tab:summary} uses $\tau^{E}=1$, a complete price freeze;
Table~\ref{tab:dose} sweeps $\tau^{E}\in[0,1]$. Under the \emph{price} shock the
adoption loss is almost exactly linear in $\tau^{E}$. Under the \emph{supply}
shock it is convex: small caps barely touch adoption and the collapse is
concentrated near the full freeze. A partial shield is therefore much less
damaging to the transition when supply is inelastic---but, as the next result
shows, much more damaging to everything else.

\input{output/tab_dose_core_import.tex}

\subsection{Welfare, and the supply-elasticity sign flip}
\label{sec:welfare}

Under the supply shock the cap yields cumulative output of $-45.6$ against
$-13.9$ with no policy: with inelastic supply, capping the price prevents
substitution away from energy and pushes the economy towards Leontief. This
survives the addition of an adoption margin.



Table~\ref{tab:summary_price_import} turns to welfare, measured as the total
consumption-equivalent variation (CEV) relative to the common no-tax steady
state. On private welfare the Slutsky transfer dominates the cap under
\emph{both} shocks (price: $-0.82\%$ vs.\ $-1.15\%$; supply: $-0.29\%$ vs.\ $-1.08\%$); the case for the transfer over the cap
does not rest on the transition. But the \emph{direction} of the cap's welfare
effect turns on the supply elasticity: tightening the cap raises welfare under
the price shock and lowers it under the supply shock (Figure~\ref{fig:dose}).

\begin{figure}[H]\centering
\includegraphics[width=\textwidth]{output/fig6_dose_response.png}
\caption{Dose-response in $\tau^{E}$ (orange) against the Slutsky transfer
benchmark (dotted). Note the sign reversal of the welfare panel between the two
shocks.}
\label{fig:dose}
\end{figure}

\input{output/tab_summary_import.tex}

\paragraph{Locating the Langot--Bayer contrast.} \citet{Langot2026},
assuming elastic supply, find the French shield effective; \citet{Bayer2026},
assuming fixed union-wide supply, find subsidies destructive. Our two closures
reproduce the \emph{sign} each obtains, which locates one source of the contrast
in a single parameter. The narrow point is that, once the supply elasticity is fixed, the direction of the
cap's welfare effect is determined---so part of the empirical disagreement reduces
to measuring that elasticity at the relevant horizon.

\paragraph{Private welfare and the transition are distinct.} The CEV above is a
\emph{private} household-welfare measure: it values consumption and leisure along
the transition and prices no climate or transition externality. Two results are
therefore kept separate. On private welfare alone the transfer dominates the cap
under both shocks; the case for the transfer does not rest on the transition.
Separately, the transfer preserves the green adoption margin the cap eliminates
(Section~\ref{sec:cap})---a positive statement about the transition, not a welfare
ranking. We do not aggregate the two into one social-welfare number: that would
require a social value of adoption (a carbon price times avoided brown-energy use)
we do not model. A positive such value only reinforces the private ranking; it
cannot reverse it, since the transfer already dominates.

\subsection{Distributional incidence and regressive targeting}
\label{sec:flat}

Without policy the impatient (high-MPC) households lose far more than the patient
(Table~\ref{tab:cev}). This makes the incidence of the transfer a first-order
design question. The Slutsky transfer indexes the lump sum to each household's
pre-crisis brown-energy use. Replacing it with an \emph{untargeted} flat transfer
of the \emph{same aggregate envelope}---the identical outlay ($49.3$ vs.\ $49.3$
over $H=24$), paid flat rather than indexed---almost halves the welfare loss, from
a CEV of $-0.82\%$ to $-0.32\%$ (Table~\ref{tab:summary_price_import}), and even
improves output ($\sum y=-18.4$ vs.\ $-22.3$), while leaving adoption intact
($9.1$ pp). The gain is entirely distributional.

The mechanism is that energy-indexed targeting is regressive \emph{in this model}.
With homothetic CES energy demand the energy share is constant across the wealth
distribution, so brown-energy use rises with total consumption, hence with
permanent income and patience. Indexing the transfer to pre-crisis energy use
therefore hands the largest cheques to the low-MPC, wealthy households---exactly
where a euro buys the least insurance---while a flat transfer reaches the high-MPC
households the crisis hits hardest. This speaks to the 2022--23 European energy
``brakes,'' which were indexed to historical consumption: within the model the
indexing is the source of the inefficiency, not a refinement of it.

We flag the caveat plainly. The result rests on the homothetic energy nest: if
energy were a necessity with a declining Engel curve---the empirically standard
case---low-income households would spend a \emph{larger} share and the sign could
reverse. The clean statement is conditional: \emph{under constant energy shares,
energy-indexed targeting is dominated by an equal-value flat transfer}.

\input{output/tab_cev_core_import.tex}
